\section{Python}

Python provides the following features:

\begin{itemize}
	\item{Positional Value}
		\begin{itemize}
			\item{Callers typically decide between naming or positioning
			      values but authors can restrict this decision.}
			\item{Positional values must precede named values.}
		\end{itemize}
	\item{Default Value}
	\item{Named Value}
	\item{Typed Value}
	\item{Caller Destructuring}
	\begin{itemize}
		\item{Restricted to iterables and dictionaries.}
	\end{itemize}
	\item{Variadic Functions}
		\begin{itemize}
		\item{One or two variadic variables will be bound to a list for
		      positional values and/or a dictionary for named values.}
		\end{itemize}
\end{itemize}

Python used to provide author destructuring for tuples, but this was removed in version
3.0.

\subsection{Background}
Python has a compiler which produces bytecode
\parencite[][internals/compiler.rs section "Abstract"]{python-developer-guide},
and an interpreter which executes the bytecode
\parencite[][internals/interpreter.rst section "Introduction"]{python-developer-guide}.
Argument and return values are given using a stack managed by the interpreter;
instructions may modify or move these values, even if this is not the primary purpose of
the instruction
\parencite[][internals/interpreter.rst section "The Evaluation Stack"]{python-developer-guide}.

There are 2 families of instructions that are used throughout the examples in this
section.
The \texttt{LOAD} family puts values on the stack from different locations depending on
the instruction.
The \texttt{CALL} family initiates a function call after the stack has been prepared.
There are also example-specific instructions which will be explained alongside the
relevant example.

\textit{Note: This section omits bookkeeping instructions that are not topically relevant.
For example, when calling a non-method function (one which is not associated with an
object instance), the interpreter pushes \texttt{NULL} onto the stack before pushing the
function.
This instruction, and similarly uninteresting instructions, are omitted for brevity.}

\subsubsection{\texttt{LOAD} family}
Instructions prefixed with \texttt{LOAD\_} retrieve a value from some location (depending
on the instruction) and put it onto the stack.
Each instruction recieves an integer which represents an index into a C-level array.
Which array is referenced depends on the instruction.
When the array contains variable names, the instruction also retrieves the value
associated with that name.
The below table explains the contents of the array that each instruction references.

\begin{tabularx}{\columnwidth}{r X}
	\texttt{LOAD\_CONST}  &
		Constant values which appear literally or implicitly in source code
		\parencite[][Doc/library/dis.rst lines 964-966]{python-source}. \\
	\texttt{LOAD\_FAST}   &
		Names of local variables which are guaranteed to be initialized.
		\parencite[][Doc/library/dis.rst lines 1253-1259,
		             Lib/inspect.py line 514
		]{python-source} \\
	\texttt{LOAD\_NAME}   &
		Names of non-local variables.
		If a local variable exists with the same name as a non-local variable then the value
		bound to the local variable will be returned.
		\parencite[][Doc/library/dis.rst lines 969-972,
		             Lib/inspect.py line 511
		]{python-source}. \\
\end{tabularx}

\subsubsection{\texttt{CALL} Family}

These instructions tell the interpreter to call a function.
This paper views this family as rooted in the plain \texttt{CALL} instruction, with all
others being variants on this core instruction.
When it needs to reference a behavior which occurs when a function is called, it examines
only the plain \texttt{CALL} instruction and assumes that other instructions behave
similarly unless the purpose of the variant is to change that specific behavior.

\begin{tabularx}{\columnwidth}{r X}
	\texttt{CALL} &
		Receives an integer indicating the number of argument values provided by the caller.
		The stack will contain the function to call followed by the argument values in
		separate stack locations.
		\parencite[][Doc/library/dis.rst lines 1398-1410,
		             Python/ceval.c lines 1314-1536
		]{python-source}. \\
\end{tabularx}

\subsection{Simple Function and Call}

\begin{pygmented}[lang=python]
def add_values(a, b):
    return a + b

add_values(1000, 1001)
\end{pygmented}

The bytecode generated for the call to \texttt{add\_values} performs 3 tasks.
First, it pushes the function onto the stack.
Next, it pushes literal values which will become associated with arguments.
Finally, it calls the function.

\begin{verbatim}
LOAD_NAME  0 (add_values)
LOAD_CONST 1 (1000)
LOAD_CONST 2 (1001)
CALL       2
\end{verbatim}

The bytecode generated for the definition is similar.
It uses \texttt{LOAD\_FAST} (instead of \texttt{LOAD\_NAME}) to refer to the variables
associated with arguments and \texttt{BINARY\_OP} (instead of \texttt{CALL}) to use the
built-in \texttt{+} operator.

\begin{verbatim}
LOAD_FAST    0 (a)
LOAD_FAST    1 (b)
BINARY_OP    0 (+)
\end{verbatim}

\subsection{Positional Value}
Provided by default.
Generally, callers can choose whether to provide values by name or position when they
make the call.
All positional values must precede all named values
\parencite[][reference/expressions.html section 6.3.4 "Calls"]{python-doc-html}.
Function authors can specify that some arguments with only receive their value by
position.
These are referred to as "positional-only arguments".

\subsubsection{Syntax \& Semantics}
Callers specify positional values by providing a comma-separated list of values.
Function authors specify positional-only arguments by listing a literal \texttt{/} after
the final positional-only argument
\parencite[][reference/compound\_stmts.html section 8.7 "Function Definitions"]{python-doc-html}:

\begin{pygmented}[lang=python]
def add_values_mixed(position, /, either):
    return position + either

# valid
add_values_mixed(1000, 1001)
add_values_mixed(1000, b=1001)

# invalid: the value for argument
# "position" cannot be given by name
# add_values_mixed(position=1000, \
#                  either=1001)
\end{pygmented}

\subsubsection{Implementation}
The bytecode for function definitions is identical to the simple definition regardless of
the presence of positional-only arguments.
The restriction is enforced within the \texttt{CALL} instruction.
In particular, the helper function \texttt{positional\_only\_passed\_as\_keyword} uses the
\texttt{co\_posonlyargcount} and \texttt{co\_localsplusnames} members of the code object.
These variables track the number of positional-only arguments
\parencite[][Doc/library/inspect.rst lines 180-181]{python-source}
and the names of all arguments
\parencite[][Include/cpython/code.h line 155]{python-source}
respectively.
If any names of positional-only arguments appear as keyword arguments then the helper
raises an error.
\parencite[][Python/ceval.c lines 1182-1244]{python-source}.
Note that the helper is only called if the function does not include a variadic variable
for named values
\parencite[][Python/ceval.c lines 1417-1431]{python-source}.

The bytecode for positional values is identical to the bytecode for the simple call.
Python stores the values of local variables in the C-level array `localsplus`.
The \texttt{CALL} instruction copies positional arguments from the stack into this array
\parencite[][Python/ceval.c lines 1341-1353]{python-source}.

\subsubsection{Historical Record}
Positional values have always been available in Python and requires no special syntax to
use.\footnote{Unfortunately, positional values are assumed to be the default method of
passing arguments by most programmers, including language authors, so there is no good
citation for this assertion.}

PEP 570 introduced positional-only arguments
\parencite[][peps/pep-0570.rst]{python-peps}.
It gives several justifications for the change, most of which are concerned with
maintaining a healthy ecosystem.
There are two relevant
\footnote{
	There are also several concerns mentioned which are specific to Python
	and/or its implementation.
	For example, it references PEP 399 which requires that pure Python code and C extensions
	have the same expressive power.
	While important to the Python community, these concerns are not relevant to this paper.
}
ecosystem harms the PEP is concerned with: inappropriate use of
value names by callers and increased maintenance burden for library authors.

Inappropriate use of value names includes using non-meaningful names, such as a math
function that takes one argument (the \texttt{sqrt} function takes one argument named
\texttt{x}).
It also includes providing values in an illogical order, such as calling the
\texttt{range} function and supplying the \texttt{stop} value before the \texttt{start}
value.

The increased maintenance burden occurs because all argument names are automatically and
irrevocably added to the API surface of all libraries.
It could be the case that a library author wants to implement a change which should be
non-breaking in principle, but prompts a variable name change for clarity.
This variable name change transforms the overall change into a breaking change.

The PEP is also concerned with functions that include a variadic variable for named
parameters.
For these functions, any non-variadic variables restrict the domain of the variadic
variable, as their names will be associated with the distinct variable rather than the
variadic one.

Finally, the PEP notes the curious case of the \texttt{range} function, which the PEP
describes as accepting "an optional parameter to the left of its required parameter."
In particular, if the \texttt{range} function only receives a single argument it is
interpreted as the end of the range, but if it receives 2 arguments then the first is
interpreted as the start while the second is interpreted as the end.
This concern does not appear to be addressed by PEP 570
\footnote{
	At least, I do not see anything that addresses it when I read the PEP and the
	implementation of range still inspects the number of provided arguments manually
	\parencite[][Objects/rangeobject.c lines 81-120]{python-source}.
}.

\subsection{Default Value}
Provided by default argument values.

\subsubsection{Syntax \& Semantics}
Function authors can define a default value by adding a literal \texttt{=} after the name
of an argument, then the value
\parencite[][reference/compound\_stmts.html section 8.7 "Function Definitions"]{python-doc-html}.

\begin{pygmented}[lang=python]
def add_values(mandatory, optional=2000):
    return mandatory + optional

# valid
add_values(1000)
add_values(1000, 1001)

# invalid: the first argument is not
# associated with a default value
# add_values()
\end{pygmented}

\subsubsection{Implementation}
Default values do not impact the bytecode generated for the definition or the call: they
are both identical to the simple call.
Instead, the \texttt{CALL} instruction retrieves default values from the code object
and uses them when necessary
\parencite[][Python/ceval.c lines 1314-1536, trace\_call\_TO\_initialize\_locals]{python-source}.

\subsubsection{Historical Record}
Default values were added in version 1.0.2
\parencite[][Misc/HISTORY lines 32809-32811]{python-source}.
There is additionally a note that default values "would now be quite sensible" in the
version 0.9.4 release notes.
This version changed argument processing so that functions receive all arguments as
separate values, rather than as a single tuple
\parencite[][Misc/HISTORY lines 34550-34639]{python-source}.

PEP 671 proposes adding a feature which would allow function authors to provide an
expression which will produce a default value at call time ("late evaluation")
\parencite[][peps/pep-0671.rst]{python-peps}.
Currently, default values must be static/constant values which are determined when the
function is defined.
The mailing list discussion includes several disagreements, including whether or not it is
appropriate for a function signature to contain un-inspectable objects and technical
difficulties about scoping rules for late evaluated values.
The proposal is still in the "draft" state, so it might be added in the future (possibly
after trivial or significant changes to the proposal), but there has been no activity on
the mailing list since 2021.
\parencite[][]{python-pep-671-mail-discussion}

\subsection{Named Value}
Provided through keyword arguments.
Generally, callers can choose whether to provide values by name or position when they
make the call.
All named values must proceed all positional values
\parencite[][reference/expressions.html section 6.3.4 "Calls"]{python-doc-html}.
Function authors can specify that some arguments will only receive their value by name.
These are referred to as "keyword-only arguments".

\subsubsection{Syntax \& Semantics}
Callers provide named values by writing first a symbolic name, then a literal \texttt{=},
then the value.
\parencite[][reference/expressions.html section 6.3.4 "Calls"]{python-doc-html}.
Function authors specify keyword-only arguments by listing a literal \texttt{*} before the
first keyword-only argument.
\parencite[][reference/compound\_stmts.html section 8.7 "Function Definitions"]{python-doc-html}.

\begin{pygmented}[lang=python]
def add_values_mixed(either, *, named):
    return either + named

# valid
add_values_mixed(named=1001, either=1000)
add_values_mixed(1000, named=1001)

# invalid: the value for argument "named"
# must be given by name
# some_function(1000, 1001)

# invalid: named values cannot appear
# before positional values
# some_function(named=1001, 1000)
\end{pygmented}

\subsubsection{Implementation}
The bytecode for function definitions is identical regardless of the presence of
keyword-only arguments.
The restriction is enforced within the \texttt{CALL} instruction.
In particular, the helper function \texttt{initialize\_locals} checks that the number of
positional arguments is not more than expected
\parencite[][Python/ceval.c lines 1458-1462, trace\_call\_TO\_initialize\_locals]{python-doc-html},
by checking the \texttt{co\_argcount} member which tracks the number of arguments which
may be positional
\parencite[][Doc/library/inspect.rst lines 146-149]{python-source}.
\parencite[][Python/assemble.c line 556]{python-source}.

Bytecode for calls which use named values differ significantly from the simple call.
For example, consider this code:

\begin{pygmented}[lang=python]
def add_values(a, b):
    return a + b

add_values(b=1001, a=1000)
add_values(1000, b=1001)
\end{pygmented}

The bytecode for the first call differs from the simple call by adding the
\texttt{KW\_NAMES} instruction prior to the \texttt{CALL} instruction:

\begin{verbatim}
LOAD_NAME                1 (add_values)
LOAD_CONST               6 (1001)
LOAD_CONST               5 (1000)
KW_NAMES                 8 (('b', 'a'))
CALL                     2
\end{verbatim}

\texttt{KW\_NAMES} marks the given constant, in this case the tuple \texttt{('b', 'a')},
as a set of argument names to be used by \texttt{CALL}
\parencite[][Python/bytecodes.c lines 2601-2605, 2644, 2869-2692, 2706-2709]{python-source}.
Then, the \texttt{CALL} instruction determines which values belong to which arguments by
corresponding their respective positions on the stack and in the tuple.
\parencite[][Python/ceval.c lines 1383-1384]{python-source}.

The process is similar when some values are provided by position and others by name.
The second call above does not have any additional instructions to handle this case:

\begin{verbatim}
LOAD_NAME                2 (add_values_mixed)
LOAD_CONST               5 (1000)
LOAD_CONST               6 (1001)
KW_NAMES                12 (('named',))
CALL                     2
\end{verbatim}

The \texttt{CALL} instruction infers which value is named based on the restriction that
positional values must precede named values
\parencite[][Python/ceval.c lines 1383-1384]{python-source}.

\subsubsection{Historical Record}
Named values were first introduced in Python 1.3
\parencite[][Doc/tut.tex lines 3540-3626]{python-source-1.3}.
The feature was based on the similar feature provided by Modula-3
\parencite[][Doc/tut.tex lines 3584-3586]{python-source-1.3}.
While keyword-only arguments (discussed later in this section) were added afterwards, the
core syntax and semantics of keyword-only arguments have remained unchanged.

PEP 3102 introduced keyword-only arguments.
It provides a single justification for the change: variadic functions cannot make use of
default values.
The PEP gives the following example:

\begin{pygmented}[lang=python]
def sortwords(*words, case_sense=False):
    pass
\end{pygmented}

If the value associated with \texttt{case\_sense} can be provided positionally then it
must be provided in every call even if the caller wants the default value of
\texttt{False}
\footnote{
	Unless the caller wants to sort the empty list. =)
}.

\subsection{Implicit Value}

This feature is not provided by Python.

\subsection{Typed Value}
Provided through type hinting.
The Python compiler and interpreter do not change their behavior based on type hints.
However, they do guarantee that the hints will be available to external tools and provide
supporting infrastructure to help the tools work correctly.
Both static analyzers and runtime checkers can make use of annotations.

\subsubsection{Syntax \& Semantics}

Type hints are specified using function annotations, as defined in PEP 3107.
This means that function authors add a colon and type name after the variable name
associated with the argument:

\begin{pygmented}[lang=python]
def typed_argument(x: str):
  pass
\end{pygmented}

\subsubsection{Implementation}

Annotations are stored as metadata in the Python object which represents the function.
Libraries can access them through the \texttt{\_\_annotations\_\_} property, which
contains a dictionary
\parencite[][peps/pep-3107.rst, "Accessing Function Annotations"]{python-peps}.

\subsubsection{Historical Record}

The foundations for type hinting were added in PEP 3102, which defines the syntax for
function annotations
\parencite[][peps/pep-3102.rst]{python-peps}.
The Python developers then waited for external community-driven tools to experiment with
different type-checking approaches.
Eventually, they took lessons learned from the community and created a set of
recommendations in PEPs 482, 483, and 484
\parencite[][peps/pep-0484.rst, "Abstract"]{python-peps}.
Much of their content addresses type theory issues, such as generics, variance, and
special types like \texttt{Any}.
Since this initial introduction there have been a number of PEPs which further clarify
best practices or provide syntactic improvements to type specifications.
However, the core mechanism that this paper is concerned with - associating a type with an
argument, regardless of how that type is specified - remains unchanged.

\subsection{Caller Destructing}
Provided through argument unpacking.
Caller destructuring is only available for iterables and mappings.

\subsubsection{Syntax \& Semantics}
This feature allows a caller to prefix one or more iterables with \texttt{*} in order to
translate their contents into a set of positional values, and/or prefix one or more
mappings with \texttt{**} to translate their contents into a set of named values.
For mappings, keys must strings naming an argument.
\parencite[][reference/expressions.html section 6.3.4 "Calls"]{python-doc-html}

\begin{pygmented}[lang=python]
def add_values(a, b, c):
    return a + b + c

# all of the below calls are equivalent
# to this:
# add_values(1000, 1001, 1002)

# destructure an iterable into positional
# values
l = [ 1000, 1001, 1002 ]
add_values(*l)

# destructure multiple iterables into
# positional values
first_part = [ 1000 ]
second_part = [ 1001, 1002 ]
add_values(*first_part, *second_part)

# destructure a mapping into named
# values
d = { 'a': 1000, 'b': 1001, 'c': 1002 }
add_values(**d)

# destructure multiple mappings into
# named values
first_part = { 'b': 1001 }
second_part = { 'a': 1000, 'c': 1002 }
add_values(**first_part, **second_part)
\end{pygmented}

\subsubsection{Implementation}
The difference between the simple call and a call which includes destructuring is best
explained by starting with the final instruction.
While the simple call uses the plain \texttt{CALL} instruction a destructuring call uses
the \texttt{CALL\_FUNCTION\_EX} instruction.
\texttt{CALL\_FUNCTION\_EX} receives either 0 or 1 which tells it whether or not there is
a mapping to destructure
\parencite[][Doc/library/dis.rst lines 1398-1410]{python-source}.

When it receives 1, there is a mapping to destructure which will be on the top of the
stack.
While the caller can use any mapping (and any number of mappings), 
\texttt{CALL\_FUNCTION\_EX} will always see a single dictionary when it executes (the
process which ensures this is discussed in more detail later in this section).
The dictionary is turned into a set of keyword arguments by interpreting the keys as names
identifying arguments.
\parencite[][Objects/call.c lines 1029-1053]{python-source}

The next item on the stack is an iterable to destructure.
In this case, \texttt{CALL\_FUNCTION\_EX} might see any iterable on the stack.
If the iterable is not a tuple it will convert it into a tuple
\parencite[][Python/bytecodes.c lines 3198-3207]{python-source}.
The elements of this tuple will be used as positional values.
\parencite[][Python/bytecodes.c line 3219]{python-source}.

When \texttt{CALL\_FUNCTION\_EX} receives 0 the process is similar, except that the top
element of the stack is an iterable and there is no mapping.

The compiler ensures that \texttt{CALL\_FUNCTION\_EX} only receives dictionaries (rather
than the arbitrary mapping object provided by the caller) with two instructions.
First, it issues a \texttt{BUILD\_MAP} instruction to place a new dictionary on the stack
\parencite[][Doc/library/dis.rst lines 1015-1023]{python-source}.
Then it adds the keys and values of each mapping object into this dictionary by repeatedly
calling the \texttt{DICT\_MERGE} instruction.
For example, this code:

\begin{pygmented}[lang=python]
add_values(*d)
\end{pygmented}

Compiles to this bytecode (note that \texttt{BUILD\_MAP} receives the value 0 to indicate
that it is building an empty dictionary):

\begin{verbatim}
LOAD_NAME         0 (add_values)
LOAD_CONST       13 (())
BUILD_MAP         0
LOAD_NAME         3 (d)
DICT_MERGE        1
CALL_FUNCTON_EX  1
\end{verbatim}

When named values are provided separately from the destructured values, the freshly
created dictionary is prepopulated with those values.
For example, this code:

\begin{pygmented}[lang=python]
d = { 'b': 1001, 'c': 1002 }
add_values(a=1000, **d)
\end{pygmented}

Compiles to this bytecode (note that in this case, \texttt{BUILD\_MAP} receives the value
\texttt{1} to indicate that there is one key-value pair on the stack):

\begin{verbatim}
LOAD_NAME         1 (add_values)
LOAD_CONST       17 (())
LOAD_CONST       11 ('a')
LOAD_CONST        3 (1000)
BUILD_MAP         1
LOAD_NAME         4 (d)
DICT_MERGE        1
CALL_FUNCTION_EX  1
\end{verbatim}

When the caller provides multiple destructured iterables, or provides literal positional
values in addition to one or more destructured iterables, the compiler issues instructions
to merge them into a list, then converts that list into a tuple.
For example, this code:

\begin{pygmented}[lang=python]
t0 = ( 1001, )
t1 = ( 1002, )
add_values(1000, *t0, *t1)
\end{pygmented}

Compiles to this bytecode:

\begin{verbatim}
LOAD_NAME        1 (add_values)
LOAD_CONST       3 (1000)
BUILD_LIST       1
LOAD_NAME        4 (t0)
LIST_EXTEND      1
LOAD_NAME        5 (t1)
LIST_EXTEND      1
CALL_INTRINSIC_1 6 (INTRINSIC_LIST_TO_TUPLE)
CALL_FUNCTION_EX 0
\end{verbatim}

If the caller provides only a single iterable to destructure, and no literal positional
values, this iterable is placed onto the stack without modification and the tuple creation
logic contained within \texttt{CALL\_FUNCTION\_EX} itself is triggered.

\subsubsection{Historical Record}
When argument unpacking was first introduced in version 1.6
\parencite[][Misc/HISTORY lines 26740-26743]{python-source},
it only allowed callers to unpack a single iterable and/or a single mapping.
For example, the call \texttt{add\_values(*first\_part, *second\_part)} would have been
illegal.
PEP 448 expanded argument unpacking so that multiple values can be destructured in the
same call \parencite[][peps/pep-0448.rst]{python-peps}.
The rationale given for this change was enhanced readability, as previously callers would
either need to build iterables/dictionaries separately or destructure them manually,
adding additional lines of code which are semantically sparse.

\subsection{Author Destructuring}
While python does not currently support authorial destructuring, it did so prior to
version 3.0 \parencite[][peps/pep-3113.rst]{python-peps}.
It allowed authors to declare that arguments should receive tuples whose values would be
bound to separate local variables:

\begin{pygmented}[lang=python]
def distance((x1, y1), (x2, y2)):
  pass
\end{pygmented}

This function would require that callers pass in 2 values which are both tuples containing
2 elements.
The values from the first tuple would be bound to local variables \texttt{x1} and
\texttt{y1}, while the values from the second would be bound to \texttt{x2} and
\texttt{y2}.

The functionality was removed through PEP 3113.
The rationale includes multiple implementation issues which are important to the Python
community but not relevant to this paper.

\subsection{Variadic Function}
Provided by arbitrary argument lists and dictionaries.
Positional values are collected by the former while named values are collected by the
latter.

\subsubsection{Syntax \& Semantics}
Function authors specify variadic-ness by specifying the name for one or two variadic
variables.
The name for the variadic list must be prefixied by a \texttt{*} while the name for the
variadic dictionary must be preceeded by a \texttt{**}
\parencite[][reference/compound\_stmts.html section 8.7 "Function Definitions"]{python-doc-html}.

\begin{pygmented}[lang=python]
from itertools import chain

def add_values(*pos_vals, **named_vals):
	return sum(chain(pos_vals, \
	                 named_vals.values()))

# All of these values appear in the
# pos_values list
add_values(1000, 1001, 1002, 1003)

# All of these values appear in the
# named_values dictionary
add_values(named_val0=1000,
           named_val1=1001,
           named_val2=1002,
           named_val3=1003)

# The values 1002 and 1003 appear in the
# pos_vals list while the names and
# values named_arg0=1000 and
# named_arg1=1001 appear in the
# named_vals dictionary
add_all_values(1002,
               1003,
               named_val0=1000,
               named_val1=1001)
\end{pygmented}

\subsubsection{Implementation}

The interpreter tracks which positional values are also variadic values by checking the
\texttt{co\_argcount} variable associated with the function's code object.
Remaining positional arguments are moved into the appropriate variadic variable, if it
exists \parencite[][Python/ceval.c lines 1355-1376]{python-source}.
It distinguishes variadic named values from non-variadic named values by checking if the
name is expected; the interpreter already has to keep track of this information because
an unrecognized value name is considered an error for non-variadic functions
\parencite[][Python/ceval.c lines 1378-1455]{python-source}.

\subsubsection{Historical Record}
PEP 468 updated the variadic variable for named values such that the author can retrieve
the syntactic order in which the values were given.
The rationale given for this change is that some users are developing APIs where order
matters, such as serialization. \parencite[][peps/pep-0468.rst]{python-peps}
